% ============================================
% Johns Hopkins Letter Template
% Assistant Professor Cover Letter – St. Olaf College (Economics, Empirical Macro)
% ============================================

\documentclass[11pt,letterpaper]{letter}

\usepackage{graphicx}
\usepackage{eso-pic}
\usepackage{geometry}
\usepackage{microtype}
\usepackage[hidelinks]{hyperref}

% ----- Page Setup -----
\geometry{left=25mm,right=25mm,top=25mm,bottom=25mm}

% ----- Logo Placement (foreground, first page only) -----
\newcommand{\PutJHULogoFG}[1]{%
  \AddToShipoutPictureFG*{%
    \AtPageUpperLeft{%
      \hspace{10mm}%
      \raisebox{-38mm}{%
        \includegraphics[width=6cm]{#1}%
      }%
    }%
  }%
}

% ----- Sender Information -----
\signature{Decory Edwards}
\address{Department of Economics \\ Johns Hopkins University \\ Baltimore, MD 21218 \\ \texttt{dedwar65@jh.edu}}

\begin{document}
\PutJHULogoFG{Images/jhulogo.png}

\begin{letter}{Search Committee \\
Department of Economics \\
St.\ Olaf College \\
1520 St.\ Olaf Avenue \\
Northfield, MN 55057, USA}

\opening{Dear Members of the Search Committee,}

I am writing to express my interest in the Assistant Professor of Economics (Empirical Macroeconomics) position at St.\ Olaf College, to begin August 2026. I am a Ph.D.\ candidate in Economics at Johns Hopkins University, advised by Professors Christopher D.\ Carroll, Nicholas W.\ Papageorge, and Stelios Fourakis, and I expect to receive my degree in May 2026. My research fields are macroeconomics, wealth and income dynamics, and computational economics.

My research agenda is unified by a focus on understanding the determinants of wealth inequality and the mechanisms through which heterogeneity in returns to wealth shapes aggregate outcomes. In my job market paper, I estimate the degree of heterogeneity in individual portfolio returns using micro data from the Survey of Consumer Finances and simulate a structural model in which households face idiosyncratic return risk. I show that allowing for persistent heterogeneity in returns substantially improves the model’s ability to match the empirical distribution of wealth, offering a tractable way to connect micro-level return dispersion to macroeconomic inequality.

In a second project, I use the Health and Retirement Study to study the relationship between interpersonal trust and portfolio performance. Building on the literature that finds a hump-shaped relationship between trust and income, I show that a similar relationship holds for individual returns to wealth measured in the HRS data, highlighting a behavioral channel through which social attitudes shape saving and investment outcomes. This work also provides new estimates of the persistent component of returns to net worth, which motivate the calibration of heterogeneity in my structural model.

St.\ Olaf’s commitment to rigorous undergraduate education in a liberal-arts setting is an excellent fit for my research and methods. The position’s emphasis on empirical macroeconomics with potential connections to financial economics, macro labor and unemployment, big/financial data, and economic history aligns closely with my agenda linking household portfolio behavior and return heterogeneity to macroeconomic outcomes. I am excited to engage students in research, advising, and data-driven inquiry that connects economic theory, credible empirical methods, and reproducible computation.

\textbf{Teaching.} My teaching experience centers on undergraduate macroeconomics and an upper-level macro policy seminar, where I have led weekly discussion sections and held regular office hours for classes ranging from 16 to 40 students. Consistent with my teaching statement, my approach is student-centered: I aim to help students ``convince themselves'' that they have moved from confusion to understanding by pairing clear scaffolding with structured office-hour coaching that builds trust and confidence. At St.\ Olaf, I am prepared to teach intermediate macroeconomic theory; introductory and advanced econometrics (including cross-sectional and time-series); principles of economics; and advanced electives in my fields. I welcome the department’s five-course annual teaching load and the chance to mentor undergraduate research and honors projects using modern, open, and replicable workflows (e.g., Python/Stata and version control).

I believe I would be a strong fit because my computational and empirical toolkit allows me to (i) produce robust, data-backed estimates of returns and distributional outcomes, (ii) build and validate models that speak to policy-relevant questions at the intersection of macroeconomics, finance, and labor, and (iii) collaborate across teams to present insights in concise, actionable formats for diverse audiences. I am enthusiastic about contributing to a collegial department that values excellent teaching, active scholarship, and undergraduate engagement, and I appreciate the support St.\ Olaf provides for research start-up and professional activity.

I would be honored to join St.\ Olaf’s Department of Economics and to mentor students pursuing quantitative, policy, and graduate-study paths. Thank you for your consideration.

\closing{Sincerely,}

\end{letter}
\end{document}
